\section{Catálogo de características efectivas}\label{sec:catalogo}
Las fórmulas de la tabla \ref{tab:caracteristicas} corresponden a Ventas y son transcripción del contrato estructurado, antes de imputación y estandarización. $q_1,q_2,q_3$ denotan unidades de los tres meses anteriores. Las unidades monetarias son las de la extracción: no se infiere una moneda empresarial sin metadatos adicionales. El catálogo financiero independiente de 16 entradas aparece en la tabla \ref{tab:tf-caracteristicas}.
\begingroup\small
\begin{longtable}{r>{\raggedright\arraybackslash}p{6.7cm}>{\raggedright\arraybackslash}p{6.2cm}}\caption{Entradas efectivas y fórmulas del contrato. $q_j=Q_{p,t-j}$.}\label{tab:caracteristicas}\\\toprule Índice & Nombre & Fórmula registrada\\\midrule\endfirsthead\toprule Índice & Nombre & Fórmula registrada\\\midrule\endhead
1 & \texttt{\detokenize{unidades_lag_1}} & Q(p,t-1) \\
2 & \texttt{\detokenize{unidades_lag_2}} & Q(p,t-2) \\
3 & \texttt{\detokenize{unidades_lag_3}} & Q(p,t-3) \\
4 & \texttt{\detokenize{unidades_media_3}} & (q1+q2+q3)/3 \\
5 & \texttt{\detokenize{unidades_mediana_3}} & mediana(q1,q2,q3) \\
6 & \texttt{\detokenize{unidades_std_3}} & std(q1,q2,q3), ddof=0 \\
7 & \texttt{\detokenize{cambio_unidades_1}} & q1-q2 \\
8 & \texttt{\detokenize{variacion_relativa_1}} & (q1-q2)/q2 si q2>0; nulo en otro caso \\
9 & \texttt{\detokenize{tendencia_unidades_3}} & (q1-q3)/2 \\
10 & \texttt{\detokenize{facturas_lag_1}} & COUNT DISTINCT InvoiceID(p,t-1) \\
11 & \texttt{\detokenize{facturas_media_3}} & media de facturas distintas en t-1,t-2,t-3 \\
12 & \texttt{\detokenize{ingresos_lag_1}} & SUM VentaSinImpuesto(p,t-1) \\
13 & \texttt{\detokenize{ingresos_media_3}} & media de ingresos de t-1,t-2,t-3 \\
14 & \texttt{\detokenize{precio_unitario_ponderado_lag_1}} & ingresos\_lag\_1/q1 si q1>0; nulo si q1=0 \\
15 & \texttt{\detokenize{clientes_distintos_lag_1}} & COUNT DISTINCT NombreCliente(p,t-1) \\
16 & \texttt{\detokenize{unidades_por_factura_lag_1}} & q1/facturas\_lag\_1 si facturas\_lag\_1>0; nulo si no \\
17 & \texttt{\detokenize{antiguedad_observada_meses}} & ordinal(t)-ordinal(primer mes observado antes del corte) \\
18 & \texttt{\detokenize{mes_objetivo_sin}} & sin(2*pi*(mes(t)-1)/12) \\
19 & \texttt{\detokenize{mes_objetivo_cos}} & cos(2*pi*(mes(t)-1)/12) \\
20 & \texttt{\detokenize{dias_mes_objetivo}} & días calendario de t, con bisiestos \\
\bottomrule\end{longtable}
\endgroup

\section{Trazabilidad y reproducción documental}
\subsection{Relación entre afirmaciones y evidencia}
El archivo \archivo{informe/evidencias/trazabilidad.json} relaciona cada tabla y figura con su fuente y hash SHA-256. Las afirmaciones cuantitativas sobre convergencia se derivan de los CSV por época y del registro de inicialización; las de evaluación se contrastan con las matrices, predicciones y análisis estadístico. Los pesos se comparan exactamente con el estado del archivo de congelación. Las figuras copiadas mantienen los mismos bytes de su origen; las nuevas vistas de soporte temporal, comparación global con cinco métricas, convergencia, paneles de error y error por partición se generan desde registros guardados.

La procedencia principal se organiza así:
\begin{itemize}
\item Calidad, grano y cardinalidades: \archivo{resultados/auditoria/}, especialmente \archivo{perfilado_datos.json} y \archivo{comparacion_csv.json}.
\item Fórmulas, parámetros y particiones: \archivo{datos/Procesados/fase_03/} y su manifiesto.
\item Pesos y recorrido del perceptrón: \archivo{resultados/fase_04/modelo_seleccionado.json} e \archivo{historial_orden_barajado.csv}.
\item Logística, gradientes y selección: \archivo{resultados/fase_05/}, sus JSON y \archivo{historial_eta_010.csv}.
\item PCA y estados congelados: \archivo{resultados/fase_06/}.
\item Métricas y predicciones finales alineadas: \archivo{resultados/fase_07/}.
\item Descomposición temporal, acuerdos, distribuciones y sensibilidad agrupada: \archivo{resultados/fase_08/}.
\end{itemize}
Los nombres técnicos de directorio se mantienen para reproducibilidad; no organizan la narrativa académica.

\subsection{Comandos sin reentrenamiento}
Desde la raíz del proyecto, el exportador documental lee artefactos y escribe únicamente dentro de \archivo{informe/}:
\begin{verbatim}
python informe/herramientas/generar_transacciones.py
cd informe
pdflatex --disable-installer -interaction=nonstopmode -halt-on-error main.tex
pdflatex --disable-installer -interaction=nonstopmode -halt-on-error main.tex
pdflatex --disable-installer -interaction=nonstopmode -halt-on-error main.tex
\end{verbatim}
No ejecutar los entrenadores, la preparación del dataset ni la evaluación final para compilar. Los archivos \archivo{README.md} y \archivo{evidencias/verificacion_informe.json} documentan los comandos efectivos, número de tablas/figuras y controles de integridad. Se comprueban referencias cruzadas, ecuaciones, archivos gráficos y dimensiones de tablas, además de la inspección visual del PDF.

Los antecedentes de Ventas registran 104 pruebas satisfactorias antes de su elaboración documental. La evaluación financiera posterior registró 313 pruebas satisfactorias, según su archivo de validaciones. Son resultados históricos, no una suite reejecutada para esta edición. El verificador documental comprueba integridad, transcripción y reproducción de tablas/figuras con operaciones experimentales bloqueadas. Las pruebas matemáticas históricas usan ejemplos de control y no deben confundirse con reentrenamientos de WWI.

\subsection{Parámetros y series que no deben confundirse}
Los JSON conservan precisión completa de pesos y sesgos; los porcentajes del texto se redondean solo para lectura. Las curvas de accuracy y BA registran estados por época; las métricas finales corresponden al checkpoint restaurado. La norma del gradiente logístico incluye el sesgo y se registra después de cada actualización; \archivo{gradient_used_norm} corresponde al estado previo. No se transforma una curva de entrenamiento en una curva de test.

\subsection{Trazabilidad de la ampliación financiera}
La exportación documental financiera se ejecuta mediante \archivo{informe/herramientas/generar_transacciones.py}; su verificación y reproducción sin operaciones de aprendizaje o predicción corresponden a \archivo{informe/herramientas/verificar_transacciones_informe.py}. La trazabilidad adicional está en \archivo{informe/evidencias/transacciones/}. Las 18 figuras financieras son copias binarias de resultados aprobados. Las 13 tablas nuevas transcriben JSON/CSV estructurados; FNR se obtiene algebraicamente de la matriz de confusión, sin recalcular predicciones.

Los comandos de compilación de este anexo siguen vigentes. Los exportadores y verificadores iniciales de Ventas se conservan como evidencia de la versión previa; sus conteos de 21 figuras y 13 tablas describen esa versión, no la ampliación actual. El README actualizado indica el procedimiento vigente. No se ejecutan preparación de datos, entrenamiento, evaluación o PCA para elaborar esta versión documental.

\section{Continuidad obligatoria del proyecto académico}
Esta versión documenta dos aplicaciones efectivamente evaluadas: Ventas y Transacciones Financieras. Las ampliaciones siguientes deberán incorporar, para cada proceso verificado, su formulación, procedencia, metodología, resultados y conclusiones parciales. Cliente o Producto no constituyen Data Marts distintos por el solo hecho de ser dimensiones. Pedidos y Compras conservan limitaciones de viabilidad; Inventario/Movimientos permanece pendiente. No se incluyen resultados vacíos o ficticios. La identificación inequívoca del ciclo actual y la evidencia operacional de cobertura, claves y unidades también siguen pendientes. El proyecto académico completo no se declara terminado.
